\chapter{Antwoorden}
\label{chap:antwoorden}

\section{Hoofdstuk 2: Haakjes wegwerken}
\subsection*{Opgave 1: Werk de haakjes weg}
\begin{multicols}{2}
\begin{enumerate}[label=\alph*)]
    \item $12x - 6y$
    \item $18x - 12y$
    \item $2x - 4y + 6$
    \item $8ax + 10ay$
    \item $4p^{2} - 5p$
    \item $-8x - 4y + 20$
    \item $a^2 + 2a - 8$
    \item $2k - k^2$
    \item $2x^{2} - yx - y^{2}$
    \item $8p^{2} - 14p + 3$
    \item $2x^{2} + xy - 6y^{2}$
\end{enumerate}
\end{multicols}

\subsection*{Opgave 2: Werk de haakjes weg door gebruik te maken van merkwaardige producten}
\begin{multicols}{2}
\begin{enumerate}[label=\alph*)]
    \item $x^{2} - 25$
    \item $x^{2} - 36$
    \item $4m^{2} - 9$
    \item $16x^{2} - 9y^{2}$
    \item $p^2 + 10p + 25$
    \item $16y^{2} - 24y + 9$
    \item $4y^{2} + 20y + 25$
    \item $4a^{2} + 24ap + 36p^{2}$
    \item $a^2 - 6a + 9$
    \item $16x^{2} - 56x + 49$
\end{enumerate}
\end{multicols}

\newpage
\section{Hoofdstuk 3: Ontbinden in factoren}
\subsection*{Opgave: ontbind in factoren}
\begin{multicols}{2}
\begin{enumerate}[label=\alph*)]
    \item $3(3x + y)$
    \item $6(a + 3)$
    \item $3(4x - 3y)$
    \item $2(2x + y - 6)$
    \item $x(x + 2)$
    \item $t(t + 12)$
    \item $(a + 4)(a - 4)$
    \item $(p + 8)(p - 8)$
    \item $(3x + 5)(3x - 5)$
    \item $(4x + 2)(4x - 2)$
    \item $(x + 5)(x + 2)$
    \item $(x + 4)(x + 2)$
    \item $(x + 7)(x - 9)$
    \item $(y + 1)(y - 6)$
    \item $(a - 8)^{2}$
    \item $(x + 6)(x - 2)$
    \item $(x + 5)^{2}$
    \item kan niet
    \item $(x + 6)^{2}$
\end{enumerate}
\end{multicols}

\newpage
\section{Hoofdstuk 4: Machten}
\subsection*{Opgave 1: Schrijf de volgende uitdrukkingen zo eenvoudig mogelijk zonder negatieve exponenten}
\begin{multicols}{2}
\begin{enumerate}[label=\alph*)]
    \item $x^7$
    \item $y^8$
    \item $5a^8$
    \item $8x^6$
    \item $y^4$
    \item $\frac{1}{m^6}$
    \item $\frac{1}{p^3}$
    \item $a$
    \item $\frac{1}{x}$
    \item $1$
    \item $y^4$
    \item $x^4$
    \item $4x^3$
    \item $p^6$
    \item $\frac{1}{b^4}$
    \item $12x^4$
    \item $8y^7$
    \item $1$
    \item $x^8$
    \item $\frac{1}{x^2}$
    \item $\frac{1}{a^2}$
    \item $3a^9$
    \item $p^4$
    \item $\frac{1}{y^5}$
    \item $\frac{1}{y^8}$
    \item $1$
\end{enumerate}
\end{multicols}

\subsection*{Opgave 2: Schrijf de volgende uitdrukkingen zo eenvoudig mogelijk zonder negatieve exponenten}
\begin{multicols}{2}
\begin{enumerate}[label=\alph*)]
    \item $x^6$
    \item $a^{10}$
    \item $8y^9$
    \item $16x^8$
    \item $a^{10}$
    \item $\frac{1}{p^8}$
    \item $\frac{1}{x^{15}}$
    \item $y^6$
    \item $16x^{12} \cdot y^8$
    \item $8x^6 \cdot y^{12}$
    \item $\frac{27}{a^6}$
    \item $\frac{25}{m^8}$
    \item $a^4 \cdot b^2 \cdot c^8$
    \item $8a^9 \cdot c^{15}$
    \item $\frac{16x^2}{y^8}$
    \item $\frac{16y^4}{x^8}$
    \item $x^4 \cdot y^2$
    \item $a^6 \cdot b^4$
    \item $\frac{b^5}{a^2}$
    \item $\frac{y^2}{x^7}$
    \item $\frac{x^2}{4}$
    \item $\frac{p^4}{3}$
    \item $\frac{4b^6}{a^3}$
    \item $\frac{9x^7}{5y^2}$
    \item $\frac{a^{10}}{b^6}$
    \item $x^{12} \cdot y^8$
\end{enumerate}
\end{multicols}


\newpage
\section{Hoofdstuk 5: Breuken}
\subsection*{Opgave 1: Vereenvoudig}
\begin{multicols}{2}
\begin{enumerate}[label=\alph*)]
    \item $\frac{1}{3xy}$
    \item $\frac{10}{x}$
    \item $\frac{4x}{3y}$
    \item $\frac{4x}{3y}$
    \item $\frac{3x}{2}$
    \item $\frac{1}{xy^2}$
    \item $\frac{4}{a^2b}$
    \item $16b$
    \item $\frac{4}{x^2}$
    \item $\frac{1}{bc}$
    \item $b^2$
    \item $\frac{1}{3y^2}$
    \item $12x^2$
    \item $\frac{1}{a^2}$
    \item $a^2$
    \item $36$
    \item $b^2c$
    \item $x^2$
\end{enumerate}
\end{multicols}

\subsection*{Opgave 2: Herleid tot 1 breuk en vereenvoudig daarna indien mogelijk}
\begin{multicols}{2}
\begin{enumerate}[label=\alph*)]
    \item $\frac{4p}{xy}$
    \item $\frac{5a}{xy}$
    \item $\frac{2k^2}{m}$
    \item $4x-2$
    \item $\frac{1}{2}$
    \item $\frac{a}{b}$
    \item $\frac{6-3x}{2y}$
    \item $\frac{3a}{x+2}$
    \item $\frac{6-3x}{6-2x}$
    \item $\frac{a-3}{a-1}$
    \item $\frac{k-3}{k+4}$
    \item $\frac{ab+2b}{ab-a}$
    \item $\frac{2x^2}{y^3}$
    \item $\frac{12p}{p-6}$
    \item $\frac{k(2k+8)}{m(5-m)} = \frac{2k^2+8k}{5m-m^2}$
    \item $c^2$
    \item $\frac{a^2}{b^2}$
    \item $\frac{20}{y}$
\end{enumerate}
\end{multicols}

\subsection*{Opgave 3: Herleid tot 1 breuk en vereenvoudig daarna indien mogelijk}
\begin{multicols}{2}
\begin{enumerate}[label=\alph*)]
    \item $\frac{-6}{x}$
    \item $\frac{1}{y}$
    \item $\frac{7a}{x+5}$
    \item $\frac{3x}{x-3}$
    \item $\frac{x^2}{x+1}$
    \item $\frac{3b-a}{b}$
    \item $\frac{b^2+a}{b}$
    \item $\frac{2x+1}{x^2-3x+2}$
    \item $\frac{5x^2+12x+2}{(x+1)(x+2)}$
    \item $-\frac{8a}{5b}$
    \item $\frac{5}{6k}$
    \item $\frac{2a^2}{a^2-b^2}$
    \item $\frac{-2y-6}{y^2-3y}$
    \item $\frac{4m}{m^2-4}$
    \item $\frac{x^2+x-1}{x^3}$
    \item $\frac{a(cd-bd+bc)}{bcd}$
    \item $\frac{a}{c^2}$
    % \item $r.$ ontbreekt
\end{enumerate}
\end{multicols}

\subsection*{Opgave 4: Vereenvoudig}
\begin{multicols}{2}
\begin{enumerate}[label=\alph*)]
    \item $\frac{2x+1}{x^2}$
    \item $2$
    \item $\frac{2b}{a+b}$
    \item $\frac{3p+2}{p^2}$
    \item $\frac{2y+x}{y-xy}$
    \item $\frac{a^2+a}{a-b}$
    \item $\frac{2x+1}{5-x}$
    \item $\frac{3k^2-1}{2k+1}$
    \item $\frac{xy+1}{xy-1}$
    \item $\frac{cd-4}{cd-1}$
    \item $\frac{a+2b}{a-2b}$
    \item $\frac{y^2}{y-x}$
    \item $\frac{ab+2b}{3ab+a}$
    \item $\frac{m^2-4km}{k^2+4km}$
    \item $\frac{2bc-4b}{c-bc}$
    \item $\frac{xy+3y}{2xy+5x}$
    \item $\frac{2d-3c}{4c+d}$
    \item $\frac{2m-5k}{3m+k}$
\end{enumerate}
\end{multicols}

\newpage
\section{Hoofdstuk 6: Lineaire vergelijkingen en functies}
\subsection*{Opgave 1: Los de volgende vergelijkingen op in $\mathbb{R}$.}
\begin{multicols}{2}
\begin{enumerate}[label=\alph*)]
    \item $x = 4{,}5$
    \item $x = 3$
    \item $x = 5$
    \item $x = -3$
    \item $p = \frac{1}{4}$
    \item $x = \frac{13}{6}$
    \item $a = \frac{3}{2}$
    \item kan niet
\end{enumerate}
\end{multicols}

\subsection*{Opgave 2: Bereken van de volgende lijnen in het vlak de richtingsco\"effici\"ent, het snijpunt met de $X$-as en het snijpunt met de $Y$-as.}
\begin{enumerate}[label=\alph*)]
    \item $\text{rc} = \frac{3}{4}, \quad S_X = \left(\frac{4}{3}, 0\right), \quad S_Y = (0, -1)$
    \item $\text{rc} = 2, \quad S_X = \left(\frac{7}{2}, 0\right), \quad S_Y = (0, -7)$
    \item $\text{rc} = 2, \quad S_X = \left(-\frac{3}{4}, 0\right), \quad S_Y = \left(0, \frac{3}{2}\right)$
    \item $\text{rc} = -\frac{1}{5}, \quad S_X = (-1, 0), \quad S_Y = \left(0, -\frac{1}{5}\right)$
    \item $\text{rc} = 0, \quad S_X = \text{nvt}, \quad S_Y = \left(0, \frac{7}{5}\right)$
    \item $\text{rc} = \frac{1}{3}, \quad S_X = (-2, 0), \quad S_Y = \left(0, \frac{2}{3}\right)$
    \item $\text{rc} = \frac{1}{2}, \quad S_X = (0, 0), \quad S_Y = (0, 0)$
    \item $\text{rc} = \frac{5}{2}, \quad S_X = \left(\frac{3}{5}, 0\right), \quad S_Y = \left(0, -\frac{3}{2}\right)$
\end{enumerate}

\subsection*{Opgave 3: Bepaal de vergelijking van de rechte lijn door het gegeven punt met de gegeven richtingsco\"effici\"ent.}
\begin{multicols}{2}
\begin{enumerate}[label=\alph*)]
    \item $y = -2x - 3$
    \item $y = -4x + 16$
    \item $y = -3$
    \item $y = 4x - 11$
\end{enumerate}
\end{multicols}

\subsection*{Opgave 4: Loodrecht of evenwijdig}
\begin{enumerate}[label=\alph*)]
    \item $a = 8$
    \item $y_1 = -x - 4$ \ en \ $y_2 = -x - 1$
    \item $l_2: y = \frac{2}{5}x + \frac{16}{5}$
    \item $l_2: y = -\frac{1}{3}x + 3$
    \item $y_2 = \frac{3}{2}x + \frac{1}{2}$
\end{enumerate}

\subsection*{Opgave 5: Snijpunten}
\begin{multicols}{2}
\begin{enumerate}[label=\alph*)]
    \item $\left(\frac{3}{2}, \frac{1}{2}\right)$
    \item $(2, 7)$
\end{enumerate}
\end{multicols}



\newpage
\section{Hoofdstuk 7: 2 vergelijkingen met 2 onbekenden}
\subsection*{Opgave 1: Los de volgende stelsels op met de substitutiemethode}
\begin{multicols}{2}
\begin{enumerate}[label=\alph*)]
    \item $\begin{cases} x = 4 \\ y = 1 \end{cases}$
    \item $\begin{cases} x = -2 \\ y = -4 \end{cases}$
    \item $\begin{cases} x = 5 \\ y = 4 \end{cases}$
    \item $\begin{cases} p = 27 \\ q = 7 \end{cases}$
    \item $\begin{cases} a = 7 \\ b = 1 \end{cases}$
    \item $\begin{cases} x = -2 \\ y = -3 \end{cases}$
\end{enumerate}
\end{multicols}

\subsection*{Opgave 2: Los de volgende stelsels op met de eliminatiemethode}
\begin{multicols}{2}
\begin{enumerate}[label=\alph*)]
    \item $\begin{cases} x = 5 \\ y = -1 \end{cases}$
    \item $\begin{cases} a = 1 \\ b = -1 \end{cases}$
    \item $\begin{cases} p = 2 \\ q = 0 \end{cases}$
    \item $\begin{cases} x = -4 \\ y = 5 \end{cases}$
    \item $\begin{cases} x = 3 \\ y = \frac{9}{2} \end{cases}$
    \item $\begin{cases} k = -2 \\ p = -3 \end{cases}$
\end{enumerate}
\end{multicols}

\subsection*{Opgave 3: Los de volgende stelsels vergelijkingen op}
\begin{multicols}{2}
\begin{enumerate}[label=\alph*)]
    \item $\begin{cases} x = 2 \\ y = 1 \end{cases}$
    \item $\begin{cases} x = -2 \\ y = 2 \end{cases}$
    \item $\begin{cases} a = 1 \\ b = 1 \end{cases}$
    \item Strijdig stelsel
    \item Strijdig stelsel
    \item $\begin{cases} n = -38 \\ t = 9 \end{cases}$
    \item $\begin{cases} x = 3 \\ y = -1 \end{cases}$
    \item $\begin{cases} x = -4 \\ y = -3 \end{cases}$
    \item $\begin{cases} x = 29 \\ y = 46 \end{cases}$
    \item Afhankelijk stelsel
    \item Afhankelijk stelsel
    \item $\begin{cases} x = -1 \\ y = 2 \end{cases}$
\end{enumerate}
\end{multicols}



\newpage
\section{Hoofdstuk 8: Goniometrie}
\subsection*{Opgave 1: Zet gegeven hoeken om naar radialen}
\begin{multicols}{2}
\begin{enumerate}[label=\alph*)]
    \item $\alpha \approx 0.49$ rad
    \item $\beta = \frac{1}{2}\pi$ rad
    \item $\gamma \approx 1.15$ rad
    \item $\alpha \approx 0.73$ rad
\end{enumerate}
\end{multicols}

\subsection*{Opgave 2: Zet gegeven hoeken om naar graden}
\begin{multicols}{2}
\begin{enumerate}[label=\alph*)]
    \item $\alpha = 30^{\circ}$
    \item $\beta = 112.5^{\circ}$
    \item $\alpha \approx 24^{\circ}$
    \item $\gamma \approx 49.27^{\circ}$
\end{enumerate}
\end{multicols}

\subsection*{Opgave 3: Rechthoekige driehoek}
\begin{enumerate}[label=\alph*)]
    \item $b \approx 16.21, \quad c \approx 24.22, \quad \beta = 42^{\circ}$
    \item $a = 5, \quad b \approx 8.66, \quad \alpha = \frac{1}{6}\pi$ rad
    \item $\alpha \approx 33.75^{\circ}, \quad b \approx 7.48$
    \item $\beta \approx 0.93$ rad, $\quad c = 10$
\end{enumerate}

\subsection*{Opgave 4: Driehoek}
\begin{enumerate}[label=\alph*)]
    \item $\beta = 34^{\circ}, \quad c \approx 9.95$
    \item $\alpha \approx 0.05$ rad, $\quad a \approx 2.12$
    \item $b \approx 26.15, \quad \beta \approx 1.01$ rad
\end{enumerate}

\subsection*{Opgave 5: Gecombineerde driehoek}
\[
    \beta \approx 39.81^{\circ}, \quad BC \approx 7.81, \quad AB \approx 8.33
\]